\documentclass[11pt]{article}

\usepackage[T1]{fontenc}
\usepackage[utf8]{inputenc}
\usepackage{lmodern}
\usepackage[margin=1in]{geometry}
\usepackage{amsmath,amssymb,amsthm,mathtools,bm}
\usepackage{booktabs,longtable,array,tabularx}
\usepackage[shortlabels]{enumitem}
\usepackage{xcolor}
\usepackage{hyperref}
\usepackage{xurl}
\usepackage{microtype}

\hypersetup{colorlinks=true,linkcolor=blue,citecolor=blue,urlcolor=blue}
\setlength{\parindent}{0pt}
\setlength{\parskip}{0.55em}
\setlength{\emergencystretch}{3em}
\renewcommand{\arraystretch}{1.12}
\newcolumntype{L}[1]{>{\raggedright\arraybackslash}p{#1}}
\newcolumntype{R}[1]{>{\raggedleft\arraybackslash}p{#1}}
\numberwithin{equation}{section}

\newtheorem{theorem}{Theorem}[section]
\newtheorem{proposition}[theorem]{Proposition}
\newtheorem{lemma}[theorem]{Lemma}
\newtheorem{definition}[theorem]{Definition}
\newtheorem{assumption}[theorem]{Assumption}
\newtheorem{corollary}[theorem]{Corollary}
\theoremstyle{remark}
\newtheorem{remark}[theorem]{Remark}

\newcommand{\OPH}{\mathrm{OPH}}
\newcommand{\chinu}{\chi_\nu}
\newcommand{\scoh}{S_{\mathrm{coh}}}
\newcommand{\nueff}{\nu_{\mathrm{eff}}}
\newcommand{\nuoph}{\nu_{\mathrm{OPH}}}
\newcommand{\gbar}{\mathbf g_b}
\newcommand{\A}{\mathcal A}
\newcommand{\I}{\mathcal I}
\newcommand{\Ralg}{\mathcal R}
\newcommand{\U}{\mathcal U}
\newcommand{\Zrec}{\mathcal Z^{\mathrm{rec}}}
\newcommand{\RecCap}{\operatorname{Cap}}
\newcommand{\Law}{\operatorname{Law}}
\newcommand{\TV}{\mathrm{TV}}
\newcommand{\dd}{\,\mathrm d}
\newcommand{\eps}{\varepsilon}
\newcommand{\rhoa}{\rho_A}
\newcommand{\Cgeom}{C_g}
\newcommand{\Gammaeff}{\Gamma_{\mathrm{eff}}}

\title{\textbf{Theoretical Bounds on \(\chinu\) in Observer-Patch Holography}\\
\large Coherent-Matter Susceptibility and the Construction Window}
\author{Bernhard Mueller and Alex Osika}
\date{Draft of May 24, 2026}
\newcommand{\OPHPaperReleaseID}{r2011}
\newcommand{\OPHPaperReleaseDate}{August 1, 2026}
\newcommand{\OPHPaperReleaseBanner}{%
\begin{center}
\small\textbf{Paper release:} \texttt{\OPHPaperReleaseID}\quad\textbf{Released:} \OPHPaperReleaseDate
\end{center}
}
\newcommand{\PaperReleaseID}{\OPHPaperReleaseID}
\newcommand{\PaperReleaseDate}{\OPHPaperReleaseDate}
\newcommand{\PaperReleaseBanner}{%
\begin{center}
\small\textbf{Paper release:} \texttt{\PaperReleaseID}\quad\textbf{Released:} \PaperReleaseDate
\end{center}
}


\begin{document}
\maketitle
\begin{center}
\small\textbf{Paper release:} \OPHPaperReleaseID\quad
\textbf{Released:} \OPHPaperReleaseDate
\end{center}

\begin{abstract}
OPH says that an ordinary tabletop gravity test on Earth should read null. The
reason is simple. The canonical dark-sector response is controlled by
\[
\rhoa=-\frac{1}{4\pi G}\nabla\cdot[(\nuoph-1)\gbar],
\]
and Earth sits deep in the high-acceleration regime where \(\nuoph\simeq1\).

The separate hypothesis is local. A coherent material state may add a
small local correction
\[
\delta\nu=\chinu\scoh
\]
to the same interpolation function. The constant is \(\chinu\). A bench
experiment measures it or bounds it.

The useful part is the definition of \(\scoh\). The scalar is built from record
capacity, record stability, predictive boundary coupling, visible overlap
coherence, stored coherent energy density, and substrate directness. It vanishes
when the record and prediction conditions fail. It is also defined on the
observer-facing quotient, so hidden implementation choices do not change it.

In the weak-field lab limit the force law is
\[
F=\frac{g^2}{4\pi G}A_\perp\,\chinu\,\Delta\scoh .
\]
For areal load \(\Sigma=M/A_\perp\), target response fraction \(f\), and
\(\Delta\scoh=\Gammaeff u/u_0\), the required value is
\[
|\chinu|_{\mathrm{req}}
=
\frac{4\pi G\Sigma f}{g}\,
\frac{u_0}{\Gammaeff u}.
\]
Write \(\Sigma_2=\Sigma/(10^2\,\mathrm{kg\,m^{-2}})\). For
\(u_0=8.5\times10^{-10}\,\mathrm{J\,m^{-3}}\), \(\Gammaeff=1\), and
\(u=10^5\)--\(10^7\,\mathrm{J\,m^{-3}}\), the calculated full-response threshold is
\[
|\chinu|_{\mathrm{req}}\simeq
f\Sigma_2\,(7.3\times10^{-23}\ \text{to}\ 7.3\times10^{-25}),
\]
scaling as \(\Gammaeff^{-1}\). Macroscopic areal loads with
\(\Sigma_2\sim1\)--\(10\) sit in the \(10^{-22}\)--\(10^{-24}\) band. The
order-unity linear ceiling over the same stored-energy range is
\(8.5\times10^{-15}\) to \(8.5\times10^{-17}\), also scaled by
\(\Gammaeff^{-1}\). For a declared response target, \(\chinu\) has to lie
between these two OPH-derived bounds. The result is a clean construction window
for the continuation, separate from the recovered OPH core.
\end{abstract}

\tableofcontents

\section{Status and Claim Boundary}

The OPH core used here consists of finite patch carriers, mismatch-lowering repair,
observer-accessible record algebras, quotient normal forms, checkpoint continuation,
support-visible BW scaling, fixed-cap generalized-entropy stationarity, and the
Jacobson-type Einstein branch. These ingredients are supplied by the consensus,
microphysics, synthesis, and compact SM/GR papers~\cite{oph-consensus,oph-micro,oph-synthesis,oph-compact}.

The dark-sector paper supplies the settled weak-field anomaly law
\begin{equation}
\rhoa
=
-\frac{1}{4\pi G}
\nabla\cdot[(\nuoph-1)\gbar],
\label{eq:canonical-dark}
\end{equation}
interpreting the anomaly as a transported modular/collar information-defect
remainder, with no new particle species postulated~\cite{oph-dark}.

The \(\chinu\) law is a continuation outside the recovered core:
\begin{equation}
\delta\nu(x,t)=\chinu\scoh(x,t).
\label{eq:chinu-law}
\end{equation}
Here \(\scoh\) is defined from OPH-native observer-facing quantities, and the
mathematical response window for \(\chinu\) is derived. No claim is made that
\(\chinu\) is nonzero.

\section{Canonical Chain}

\subsection{Fixed-cutoff carrier}

At fixed cutoff, an observer-facing OPH implementation surface is a federated
patch carrier
\begin{equation}
\mathfrak F
=
\bigl(V,E,\{\A_i\}_{i\in V},\{\I_e\}_{e\in E},
\{\pi_{i,e}\},\{\Ralg_i\},\{\U_i\}\bigr).
\end{equation}
Here \(\A_i\) are local finite algebras, \(\I_e\) are overlap interfaces,
\(\pi_{i,e}\) are visible restriction maps, \(\Ralg_i\subseteq Z(\A_i)\) are
record algebras, and \(\U_i\) are local update and repair interfaces. Physical
claims are made through visible restrictions, record algebras, and quotient-local
observables, without choosing hidden microscopic representatives~\cite{oph-micro}.

\subsection{Mismatch and repair}

The consensus paper gives a nonnegative visible mismatch functional \(\Phi\).
Accepted repairs lower the relevant mismatch on finite state spaces. Under the
declared local-diamond and repair-completeness hypotheses, the terminal
observer-facing normal form is unique on the physical quotient and independent
of asynchronous update schedule~\cite{oph-consensus}.

\subsection{Records and checkpoints}

For an observer-supporting subfederation \(U\), the microphysics paper uses a
checkpoint of the form
\begin{equation}
\mathrm{Chk}_U(t)=
\bigl(
\Ralg_U(t),
\rho_U^{\mathrm{acc}}(t),
\mathfrak I_U^{\mathrm{ext}}(t),
\nu_{\ge t},
\mathfrak B_U(t)
\bigr).
\end{equation}
If two checkpoints agree on the accessible record algebra, accessible state,
external interface, boundary-update schedule class, and provenance bundle, they
induce the same continuation law on the observer-accessible event algebra. Approximate
agreement gives a controlled total-variation bound~\cite{oph-micro}.

\subsection{Gravity and dark response}

The compact SM/GR paper keeps the recovered core and downstream continuations
separate. On the stated support-visible scaling branch, the OPH stack gives a
Jacobson-type Einstein relation; dark-sector response laws are phenomenological
continuations and do not belong to that recovered-core tier~\cite{oph-compact}.
The dark paper then models the anomaly by Eq.~\eqref{eq:canonical-dark}.

On Earth, \(g_b/a_0\sim 10^{11}\), so the canonical \(\nuoph\) term is pinned
near unity. A terrestrial tabletop anomaly attributed only to the canonical OPH
dark sector should read null. The \(\chinu\) continuation is tested against
that null baseline as a separate effect.

\section{The Coherence Scalar}

\begin{definition}[Patch coherence factors]
Let \(U\) be an observer-supporting patch subfederation at cycle \(t\), with
observer-accessible record algebra \(\Zrec_U(t)\). Let \(Y_U(t)\) be the induced
record outcome variable, \(X_{\partial U,t+1:t+h}\) the boundary packet
over horizon \(h\), and \(\Phi_U(t)\) the visible mismatch on the relevant
support. Define
\begin{align}
R_U(t;h)
&=
\left[
1-
d_{\TV}\!\left(
\Law(Y_U(t+h)),\Law(Y_U(t))
\right)
\right]e^{-\delta_{\mathrm{rec}}(t)},
\label{eq:R-factor}\\
P_U(t;h)
&=
\frac{
I\!\left(Y_U(t);X_{\partial U,t+1:t+h}\right)
}{
H\!\left(X_{\partial U,t+1:t+h}\right)+\epsilon
},
\label{eq:P-factor}\\
C_U(t)
&=
\left[
1-\frac{\Phi_U(t)}{\Phi_U^{\max}+\epsilon}
\right]_+.
\label{eq:C-factor}
\end{align}
Fix a reference capacity \(\RecCap_0>0\) and set
\begin{equation}
\widehat{\mathrm{Cap}}_U(t)=
\frac{\RecCap(\Zrec_U(t))}{\RecCap_0}.
\end{equation}
The minimal patch-coherence factor is
\begin{equation}
\mathcal K_U(t;h)=
\mathbf 1_{\mathrm{self\mbox{-}read}}\,
\widehat{\mathrm{Cap}}_U(t)\,
R_U(t;h)\,
P_U(t;h)\,
C_U(t).
\label{eq:K-factor}
\end{equation}
\end{definition}

\begin{definition}[OPH-native coherent-matter scalar]
Let \(u_{\mathrm{stored}}(U,t)\) be the stored coherent matter energy density in
the substrate mode supported by \(U\), let \(u_0\) be a fixed reference energy
density, and let \(\eps_{\mathrm{sub}}(U)\in[0,1]\) be the substrate directness
factor. Define
\begin{equation}
\scoh(U,t;h)
=
\eps_{\mathrm{sub}}(U)
\frac{u_{\mathrm{stored}}(U,t)}{u_0}
\mathcal K_U(t;h).
\label{eq:Scoh-patch}
\end{equation}
Given a partition of unity \(w_U(x)\) over observer-supporting subfederations,
define the coarse-grained field
\begin{equation}
\scoh(x,t)=\sum_U w_U(x)\scoh(U,t;h),
\qquad
\sum_U w_U(x)=1.
\label{eq:Scoh-field}
\end{equation}
\end{definition}

The continuation axiom is the single statement
\begin{equation}
\boxed{
\nueff(x,t)=\nuoph(x,t)+\chinu\scoh(x,t).
}
\label{eq:continuation-axiom}
\end{equation}

\section{Structural Properties}

\begin{proposition}[Observer-facing quotient well-definedness]
Assume two microscopic representatives induce the same observer-accessible
continuation law on \(U\), preserve the same visible mismatch score, and expose the
same stored coherent energy density and substrate class. Then they give the same
\(\scoh(U,t;h)\).
\end{proposition}

\begin{proof}
The factors \(\RecCap(\Zrec_U)\), \(R_U\), \(P_U\), and \(C_U\) are defined from
observer-accessible records, boundary-law variables, and visible mismatch
data. These are precisely the quantities that survive the observer-facing
quotient. Multiplying them by \(u_{\mathrm{stored}}\) and
\(\eps_{\mathrm{sub}}\), which are declared observable substrate data for the
continuation, does not introduce hidden representative dependence.
\end{proof}

\begin{proposition}[Null conditions]
If \(U\) has no self-read, no stable record algebra, or no predictive boundary
coupling at horizon \(h\), then \(\scoh(U,t;h)=0\).
\end{proposition}

\begin{proof}
If there is no self-read, then
\(\mathbf 1_{\mathrm{self\mbox{-}read}}=0\). If no stable record algebra is
available, then either \(\RecCap(\Zrec_U)=0\) or the stability factor \(R_U\)
vanishes on the declared readout. If predictive boundary coupling is absent,
then \(P_U=0\). In each case \(\mathcal K_U=0\), hence \(\scoh=0\).
\end{proof}

\begin{proposition}[Saturated-limit reduction]
If
\[
\mathbf 1_{\mathrm{self\mbox{-}read}}\simeq1,\quad
\widehat{\mathrm{Cap}}_U\simeq1,\quad
R_U\simeq1,\quad
P_U\simeq1,\quad
C_U\simeq1,
\]
then
\[
\scoh(U,t;h)\simeq
\eps_{\mathrm{sub}}(U)\frac{u_{\mathrm{stored}}(U,t)}{u_0}.
\]
\end{proposition}

\begin{proof}
Substitute the saturated factor values into Eq.~\eqref{eq:Scoh-patch}.
\end{proof}

\section{Weak-Field Response}

Insert Eq.~\eqref{eq:continuation-axiom} into the settled weak-field density:
\begin{equation}
\rhoa^{\mathrm{eff}}
=
-\frac{1}{4\pi G}
\nabla\cdot
\Bigl[
\bigl(\nuoph-1+\chinu\scoh\bigr)\gbar
\Bigr].
\label{eq:rho-eff}
\end{equation}
In a small terrestrial laboratory region with \(\gbar\simeq -g\,\hat z\) and
canonical \(\nuoph-1\) negligible, this becomes
\begin{equation}
\rhoa^{\mathrm{lab}}
\simeq
-\frac{1}{4\pi G}
\nabla\cdot
\bigl[\chinu\scoh\,\mathbf g\bigr].
\label{eq:rho-lab}
\end{equation}

\begin{theorem}[Planar response law]
Let a thin device have horizontal projected area \(A_\perp\), and define the
bottom-to-top coherence contrast
\[
\Delta\scoh={\scoh}_{\mathrm{bottom}}-{\scoh}_{\mathrm{top}}.
\]
Then, to leading order in the thin-device approximation, the apparent vertical
force generated by the \(\chinu\) continuation is
\begin{equation}
\boxed{
F_\chi
=
\frac{g^2}{4\pi G}A_\perp\,\chinu\,\Delta\scoh .
}
\label{eq:force-law}
\end{equation}
\end{theorem}

\begin{proof}
With \(\mathbf g=-g\hat z\), Eq.~\eqref{eq:rho-lab} gives
\[
\rhoa^{\mathrm{lab}}
=
\frac{g\chinu}{4\pi G}\frac{\partial\scoh}{\partial z}
\]
up to the sign convention for \(z\). Integrating through a thin column gives
\[
M_A=
\int\rhoa^{\mathrm{lab}}\dd V
=
-\frac{g\chinu}{4\pi G}A_\perp\,\Delta\scoh
\]
when \(\Delta\scoh\) is bottom minus top. The apparent weight change is
\(F_\chi=-gM_A\), yielding Eq.~\eqref{eq:force-law}.
\end{proof}

\begin{corollary}[Geometric ceiling]
If \(|\Delta\nu|\lesssim 1\), then
\begin{equation}
\frac{|F_\chi|}{A_\perp}
\lesssim
\Cgeom
\equiv
\frac{g^2}{4\pi G}
=
1.1466\times10^{11}\,\mathrm{N\,m^{-2}}
\label{eq:geometric-ceiling}
\end{equation}
for \(g=9.80665\,\mathrm{m\,s^{-2}}\) and
\(G=6.67430\times10^{-11}\,\mathrm{m^3\,kg^{-1}\,s^{-2}}\).
\end{corollary}

This ceiling is algebraic. It does not assert that any material system reaches
\(|\Delta\nu|=1\).

\section{The Susceptibility Window}

For bounds it is useful to package the contrast geometry, substrate directness,
and OPH coherence factors into one dimensionless number:
\begin{equation}
\Delta\scoh
=
\Gammaeff\frac{u}{u_0},
\qquad
0\le \Gammaeff\le 1.
\label{eq:Gamma-definition}
\end{equation}
Here \(u\) is the representative stored coherent energy density in the active
mode. The factor \(\Gammaeff\) includes vertical contrast, substrate directness,
and effective \(\mathcal K\).

\begin{theorem}[Nonempty macroscopic linear-response window]
Let \(\Sigma=M/A_\perp\) be the areal load, let \(f\in[0,1]\) be the target
fraction of the ordinary weight \(Mg\), and let \(\epsilon_{\mathrm{lin}}\le 1\)
be the maximum allowed \(|\Delta\nu|\) in the chosen linear-response regime.
Assume \(\Gammaeff>0\). If
\begin{equation}
f\,\frac{4\pi G\Sigma}{g}
<
\epsilon_{\mathrm{lin}},
\label{eq:window-condition}
\end{equation}
then the interval
\begin{equation}
\boxed{
\frac{4\pi G\Sigma f}{g}\,
\frac{u_0}{\Gammaeff u}
\le
|\chinu|
\le
\epsilon_{\mathrm{lin}}
\frac{u_0}{\Gammaeff u}
}
\label{eq:chinu-window}
\end{equation}
is nonempty. Values in this interval produce at least fraction \(f\) of the
ordinary weight inside the declared linear-response bound. Conversely, any
value that produces at least fraction \(f\) while staying inside that
linear-response bound lies in this interval.
\end{theorem}

\begin{proof}
The force law gives
\[
\frac{|F_\chi|}{Mg}
=
\frac{g}{4\pi G\Sigma}|\chinu|\Delta\scoh.
\]
Requiring \(|F_\chi|/(Mg)\ge f\) gives the lower bound in
Eq.~\eqref{eq:chinu-window}. Requiring
\(|\Delta\nu|=|\chinu|\Delta\scoh\le\epsilon_{\mathrm{lin}}\) gives the upper
bound. The interval is nonempty exactly when the lower bound is smaller than the
upper bound, which is Eq.~\eqref{eq:window-condition}.
\end{proof}

\begin{corollary}[Numerical width in normalized areal load]
Let
\[
\Sigma_2=\frac{\Sigma}{10^2\,\mathrm{kg\,m^{-2}}}.
\]
Then
\[
\frac{4\pi G\Sigma}{g}=8.5525\times10^{-9}\,\Sigma_2.
\]
The full-response \((f=1)\), order-unity linear window has width
\[
\frac{\chi_{\mathrm{upper}}}{\chi_{\mathrm{lower}}}
=
1.1692\times10^8\,\Sigma_2^{-1},
\]
independent of \(u\) and \(\Gammaeff\). With
\(\epsilon_{\mathrm{lin}}=0.1\), the window width is
\(1.1692\times10^7\,\Sigma_2^{-1}\).
\end{corollary}

\begin{table}[h]
\centering
\begin{tabular}{@{}R{0.19\textwidth}R{0.23\textwidth}R{0.25\textwidth}R{0.20\textwidth}@{}}
\toprule
\(u\,[\mathrm{J\,m^{-3}}]\) &
\(u/u_0\) &
\(|\chinu|_{\mathrm{req}}\) for \(f\Sigma_2/\Gammaeff=1\) &
\(|\chinu|_{\mathrm{lin}}\) for \(\epsilon_{\mathrm{lin}}/\Gammaeff=1\) \\
\midrule
\(1.0\times10^5\) & \(1.18\times10^{14}\) & \(7.27\times10^{-23}\) & \(8.5\times10^{-15}\) \\
\(5.5\times10^5\) & \(6.47\times10^{14}\) & \(1.32\times10^{-23}\) & \(1.55\times10^{-15}\) \\
\(1.0\times10^6\) & \(1.18\times10^{15}\) & \(7.27\times10^{-24}\) & \(8.5\times10^{-16}\) \\
\(4.0\times10^6\) & \(4.71\times10^{15}\) & \(1.82\times10^{-24}\) & \(2.13\times10^{-16}\) \\
\(1.0\times10^7\) & \(1.18\times10^{16}\) & \(7.27\times10^{-25}\) & \(8.5\times10^{-17}\) \\
\bottomrule
\end{tabular}
\caption{Full-response susceptibility bounds for \(u_0=8.5\times10^{-10}\,\mathrm{J\,m^{-3}}\) and \(\Sigma_2=\Sigma/(10^2\,\mathrm{kg\,m^{-2}})\). The third column should be multiplied by \(f\Sigma_2/\Gammaeff\). The fourth column should be multiplied by \(\epsilon_{\mathrm{lin}}/\Gammaeff\).}
\label{tab:window}
\end{table}

The table gives the bound in numerical form. It says that
ordinary material stress and coherent phonon energy densities put the
macroscopic response threshold in the \(\chinu\sim10^{-22}\)--\(10^{-24}\)
range for macroscopic areal loads and \(\Gammaeff\) near unity. Smaller
\(\Gammaeff\) shifts the required value upward by \(1/\Gammaeff\).

\section{Power and Stored-Energy Bounds}

The force law fixes stored energy. Input power depends on how the coherent state
is maintained.

\subsection{Driven coherent oscillations}

For a driven oscillator with angular frequency \(\omega\), mode volume
\(V_{\mathrm{mode}}\), quality factor \(Q\), and stored density \(u\),
\begin{equation}
P_{\mathrm{loss}}
=
\frac{uV_{\mathrm{mode}}\omega}{Q}.
\label{eq:loss-power}
\end{equation}
Combining Eq.~\eqref{eq:loss-power} with the required density from the
macroscopic response law gives
\begin{equation}
\boxed{
P_{\mathrm{loss}}
=
\frac{4\pi G\Sigma f\,u_0\,V_{\mathrm{mode}}\omega}
{g\,Q\,\Gammaeff\,|\chinu|}.
}
\label{eq:power-bound}
\end{equation}

For \(\Sigma_2=1\),
\(V_{\mathrm{mode}}=10^{-3}\,\mathrm{m^3}\),
\(\omega=2\pi(4.0\times10^4)\,\mathrm{s^{-1}}\),
\(Q=10^5\), \(f=1\), and \(\Gammaeff=1\), this gives:
\[
\begin{array}{c|c}
|\chinu| & P_{\mathrm{loss}}\\
\hline
10^{-20} & 1.8\,\mathrm{W}\\
10^{-21} & 18\,\mathrm{W}\\
10^{-22} & 183\,\mathrm{W}\\
10^{-23} & 1.8\,\mathrm{kW}\\
10^{-24} & 18\,\mathrm{kW}
\end{array}
\]
The scaling is \(P_{\mathrm{loss}}\propto(Q\Gammaeff|\chinu|)^{-1}\).

\subsection{Static elastic storage}

For a static elastic-strain coherent variable, the stored density is
\begin{equation}
u_{\mathrm{elastic}}=\frac{\sigma^2}{2E},
\label{eq:elastic-density}
\end{equation}
where \(\sigma\) is stress and \(E\) is Young's modulus. In that interpretation
there is no oscillator decay factor \(\omega/Q\). Maintaining the state is then
limited by creep, relaxation, and mechanical failure. Its loss law differs from
Eq.~\eqref{eq:loss-power}. The same susceptibility bounds apply, because they
depend only on \(u\), \(\Gammaeff\), and \(\Sigma\).

\section{Construction-Scale Consequences}

The bounds are algebraic. They do not prove that any substrate realizes
\(\Gammaeff\simeq1\), that \(\chinu\ne0\), or that a high-energy coherent state
can be maintained safely. They specify the requirements for a macroscopic
force device if the continuation is real.

A force device, in this narrow sense, is a bounded substrate whose active region
creates a vertical coherence contrast \(\Delta\scoh\). The apparent weight
fraction is
\begin{equation}
\frac{|F_\chi|}{Mg}
=
\frac{g}{4\pi G\Sigma}\,
|\chinu|\,\Gammaeff\frac{u}{u_0}.
\label{eq:force-fraction}
\end{equation}
All construction choices enter through the product
\begin{equation}
|\chinu|\,\Gammaeff\,u.
\label{eq:construction-product}
\end{equation}
This is why \(\chinu\) is the governing number. Geometry fixes the coefficient
\(g/(4\pi G\Sigma)\). Better materials can raise \(u\). Better mode design can
raise \(\Gammaeff\). A smaller \(\chinu\) must be paid for linearly by higher
stored energy density, larger active volume, better quality factor, or a smaller
target response fraction. If \(\chinu=0\), no amount of construction produces the
continuation force.

For the normalized areal load \(\Sigma_2=1\), the full-response condition is
\begin{equation}
\Delta\nu_{\mathrm{full}}(\Sigma_2=1)
=
\frac{4\pi G\Sigma}{g}
=
8.5525\times10^{-9}.
\label{eq:full-response-dnu}
\end{equation}
This is far below even a conservative \(\epsilon_{\mathrm{lin}}=0.1\) linear
bound. The required \(\Delta\nu\) lies inside the theorem window. Construction
depends on the required \(\Gammaeff u\) product, mode volume, quality factor,
and material stress.

\begin{table}[p]
\centering
\small
\begin{tabularx}{\textwidth}{@{}R{0.12\textwidth}R{0.16\textwidth}R{0.16\textwidth}L{0.19\textwidth}X@{}}
\toprule
\(|\chinu|\) &
\(u_{\mathrm{req}}\) at \(f\Sigma_2/\Gammaeff=1\) &
\(P_{\mathrm{loss}}\) at \(f\Sigma_2/\Gammaeff=1,Q=10^5\) &
Construction class &
Interpretation \\
\midrule
\(10^{-20}\) &
\(7.3\times10^2\,\mathrm{J\,m^{-3}}\) &
\(1.8\,\mathrm{W}\) &
small bench article &
The force is easy to produce if the coherence scalar is real. Existing ambient-coherence nulls become the main consistency check. \\
\(10^{-21}\) &
\(7.3\times10^3\,\mathrm{J\,m^{-3}}\) &
\(18\,\mathrm{W}\) &
tabletop module &
Accessible to modest driven resonators or static strain specimens. A direct balance or pendulum experiment should see it cleanly. \\
\(10^{-22}\) &
\(7.3\times10^4\,\mathrm{J\,m^{-3}}\) &
\(183\,\mathrm{W}\) &
high-Q insert / compact platform &
This is the first useful construction band. Full response is compatible with ordinary high-Q crystal or metal inserts. Thermal management matters. \\
\(10^{-23}\) &
\(7.3\times10^5\,\mathrm{J\,m^{-3}}\) &
\(1.8\,\mathrm{kW}\) &
engineered crystal or metal module &
Inside high-strength material energy densities. Continuous driven operation becomes hard. Static elastic storage becomes much more attractive. \\
\(10^{-24}\) &
\(7.3\times10^6\,\mathrm{J\,m^{-3}}\) &
\(18\,\mathrm{kW}\) &
industrial multi-module system &
Near the upper end of practical stored-density claims. Driven operation needs large area, high \(Q\), cryogenic or exceptional materials, or partial-lift use. \\
\(10^{-25}\) &
\(7.3\times10^7\,\mathrm{J\,m^{-3}}\) &
\(183\,\mathrm{kW}\) &
large array / static-only candidate &
Inside the response window. Compact driven construction is implausible. Static prestress, very high \(\Gammaeff\), or large arrays are the plausible routes. \\
\(10^{-26}\) &
\(7.3\times10^8\,\mathrm{J\,m^{-3}}\) &
\(1.8\,\mathrm{MW}\) &
outside compact construction &
The response inequalities are satisfied. Material and power requirements exceed the intended compact-device scale. This is effectively a null for compact devices. \\
\bottomrule
\end{tabularx}
\caption{Construction regimes inside the response-window theorem for \(\Sigma_2=1\), \(f=1\), \(u_0=8.5\times10^{-10}\,\mathrm{J\,m^{-3}}\), \(V_{\mathrm{mode}}=10^{-3}\,\mathrm{m^3}\), \(\omega=2\pi(4.0\times10^4)\,\mathrm{s^{-1}}\), and \(Q=10^5\). Every row has \(\Delta\nu=8.5525\times10^{-9}\), so construction difficulty changes while algebraic admissibility is the same. Multiply the listed material and power requirements by \(f\Sigma_2/\Gammaeff\).}
\label{tab:construction-regimes}
\end{table}

Two consequences follow.

First, a measurement of \(\chinu\) selects a construction class. If
\(|\chinu|\gtrsim10^{-22}\) and \(\Gammaeff\) is reasonably large, compact
demonstrators are in the ordinary resonator and high-strength-substrate regime.
If \(|\chinu|\sim10^{-23}\)--\(10^{-24}\), the design must move toward high-Q
crystals, larger active area, static strain, or partial lift. If
\(|\chinu|\lesssim10^{-25}\), compact force-device construction stops being the
right first target. The useful output is a bound on the continuation.

Second, the static-versus-driven distinction is decisive. The driven column pays
\(uV_{\mathrm{mode}}\omega/Q\). Static elastic storage pays no oscillator
maintenance cost. It must survive creep and fracture. The same
\(\chinu\) value can be impractical for a driven phonon design and practical for
a static strain design. The static-stress experiment decides which maintenance
law applies to the construction table.

\section{Experimental Bound Formula}

Let an experiment have force sensitivity \(F_{\min}\), projected area
\(A_\perp\), and known coherence contrast \(\Delta\scoh\). A null result gives
\begin{equation}
\boxed{
|\chinu|
\le
\frac{4\pi G F_{\min}}{g^2 A_\perp|\Delta\scoh|}.
}
\label{eq:null-bound-general}
\end{equation}
If the mode is a driven oscillator, with
\(\Delta\scoh=\Gammaeff QP_{\mathrm{in}}/(u_0\omega V_{\mathrm{mode}})\), then
\begin{equation}
\boxed{
|\chinu|
\le
\frac{4\pi G F_{\min}u_0\omega V_{\mathrm{mode}}}
{g^2 A_\perp\Gammaeff QP_{\mathrm{in}}}.
}
\label{eq:null-bound-driven}
\end{equation}
If the mode is static elastic storage, replace the driven density by
\(u_{\mathrm{elastic}}\) from Eq.~\eqref{eq:elastic-density}.

\section{Compatibility with Existing Nulls}

\begin{proposition}[No ordinary-matter anomaly from null coherence]
If ordinary non-engineered configurations have \(\mathcal K_U=0\) on the
observer-facing scalar of Eq.~\eqref{eq:K-factor}, then they impose no direct
bound on \(\chinu\), because \(\delta\nu=\chinu\scoh=0\) for those
configurations.
\end{proposition}

\begin{proof}
This is immediate from Eq.~\eqref{eq:Scoh-patch}. The continuation couples only
to the OPH-native coherent scalar. Undifferentiated rest mass or heat is outside
this ansatz.
\end{proof}

If ordinary matter has a small nonzero ambient scalar
\({\scoh}_{\mathrm{amb}}\), then existing gravitational null tests bound the
product:
\begin{equation}
|\chinu|\,{\scoh}_{\mathrm{amb}}
\le
\delta\nu_{\mathrm{test}}.
\label{eq:ambient-bound}
\end{equation}
This is a real consistency condition. It becomes a theorem-grade number when
the ambient scalar is specified. The continuation is falsifiable in
two ways: by direct \(\chinu\) searches through Eq.~\eqref{eq:null-bound-general}
and by ordinary-gravity null tests once a model for \({\scoh}_{\mathrm{amb}}\) is
fixed.

\section{What OPH Alone Bounds}

The OPH-only result is conditional. For a declared areal load \(\Sigma\), target
fraction \(f\), stored coherent density \(u\), effective coherence
\(\Gammaeff\), and linear allowance \(\epsilon_{\mathrm{lin}}\), any
\(\chinu\) that produces the target response while staying inside the linear
branch has to satisfy Eq.~\eqref{eq:chinu-window}. Below the lower edge, the
response is smaller than the target. Above the upper edge, the declared
linear-response branch stops applying.

\begin{proposition}[No OPH-only point value]
The recovered OPH core, together with the continuation axiom
\(\delta\nu=\chinu\scoh\), does not fix a nonzero value of \(\chinu\), its sign,
or a narrower absolute numerical interval than Eq.~\eqref{eq:chinu-window}.
\end{proposition}

\begin{proof}
The continuation term is outside the recovered Einstein/dark-sector core. The
choice \(\chinu=0\) returns exactly to that core and satisfies the quotient,
checkpoint, null-coherence, and terrestrial high-acceleration baseline
statements. The OPH construction also defines the response through the product
\(\chinu\scoh\). Without a further theorem fixing a global maximum for
\(\scoh\), a canonical material value of \(u\), or a variational principle for
the continuation coefficient, the core theorems constrain only products and
conditional response windows. Finite patch capacity makes the scalar
well-defined at fixed cutoff; it does not select the numerical coupling.
\end{proof}

Further narrowing needs one extra input: a canonical normalization of
the maximum ordinary coherent scalar, an ambient-coherence theorem for ordinary
matter, or a direct experimental bound. Those inputs can sharpen
Eq.~\eqref{eq:chinu-window}. They are separate from the recovered OPH core.

\section{Open Normalizations}

Three normalizations are parameters of the continuation.

First, \(\RecCap_0\) must be fixed canonically. Depending on the substrate and
regime, \(\RecCap\) could be \(\log\dim\Zrec\), an effective rank, usable record
entropy, or effective patch capacity. The formulas keep \(\RecCap_0\) explicit.

Second, the partition weights \(w_U(x)\) should be induced by a declared
observer-supporting subcover. Only the coarse-grained contrast
\(\Delta\scoh\) enters the bounds, so the construction is an implementation
choice.

Third, the exponents in Eq.~\eqref{eq:K-factor} are all set to one. A calibrated
variant may replace
\[
\mathcal K_U=
\mathbf 1_{\mathrm{self\mbox{-}read}}\,
\widehat{\mathrm{Cap}}_U R_U P_U C_U
\]
by a renormalized exponent family. The linear product is the minimal
non-arbitrary continuation.

\section{Conclusion}

The \(\chinu\) continuation can be stated as one additional local response law:
\[
\nueff=\nuoph+\chinu\scoh.
\]
The scalar \(\scoh\) can be built from OPH-native observer-facing quantities:
record capacity, record stability, boundary predictivity, visible overlap
coherence, and stored coherent matter energy density. It is null when the
observer-facing coherence conditions fail and quotient-invariant when the
observer-accessible continuation law is unchanged.

The response bounds are simple. For an areal load \(\Sigma\), target fraction
\(f\), and contrast \(\Delta\scoh=\Gammaeff u/u_0\),
\[
|\chinu|_{\mathrm{req}}
=
\frac{4\pi G\Sigma f}{g}\frac{u_0}{\Gammaeff u}.
\]
The corresponding linear headroom extends up to
\[
|\chinu|_{\mathrm{lin}}
\le
\epsilon_{\mathrm{lin}}\frac{u_0}{\Gammaeff u}.
\]
For \(\Sigma_2=\Sigma/(10^2\,\mathrm{kg\,m^{-2}})\), the interval width is
\(1.1692\times10^8\,\epsilon_{\mathrm{lin}}/(f\Sigma_2)\). Coherent stored
energy densities \(10^5\)--\(10^7\,\mathrm{J\,m^{-3}}\) put macroscopic
full-response thresholds near \(10^{-22}\)--\(10^{-24}\), scaled by
\(f\Sigma_2\Gammaeff^{-1}\). This is a mathematical continuation window outside
the recovered OPH predictions and outside experimental results. The recovered
core permits \(\chinu=0\); a nonzero value or a tighter interval requires the
additional inputs listed above.

\begin{thebibliography}{99}

\bibitem{oph-synthesis}
B.\ M\"uller, A.\ Osika, K.\ Xue, B.\ Cassie, P.\ Nguyen, M.\ Poneder, and K.\ A.\ Anirudha,
\emph{Observers Are All You Need}.
\href{https://github.com/FloatingPragma/observer-patch-holography/releases/download/\OPHPaperReleaseID/observers_are_all_you_need.pdf}{release PDF}.

\bibitem{oph-compact}
B.\ M\"uller, A.\ Osika, K.\ Xue, and P.\ Nguyen,
\emph{Recovering Relativity and Standard Model Structure from Observer-Overlap Consistency}.
\href{https://github.com/FloatingPragma/observer-patch-holography/releases/download/\OPHPaperReleaseID/recovering_relativity_and_standard_model_structure_from_observer_overlap_consistency_compact.pdf}{release PDF}.

\bibitem{oph-consensus}
B.\ M\"uller, A.\ Osika, K.\ Xue, and P.\ Nguyen,
\emph{Reality as a Consensus Protocol}.
\href{https://github.com/FloatingPragma/observer-patch-holography/releases/download/\OPHPaperReleaseID/reality_as_consensus_protocol.pdf}{release PDF}.

\bibitem{oph-micro}
B.\ M\"uller and K.\ Xue,
\emph{Federated Echosahedral Screen Microphysics: Patch Hardware, Records, and Observer Synchronization in OPH}.
\href{https://github.com/FloatingPragma/observer-patch-holography/releases/download/\OPHPaperReleaseID/screen_microphysics_and_observer_synchronization.pdf}{release PDF}.

\bibitem{oph-calibration}
B.\ M\"uller and K.\ Xue,
\emph{Digital Calibration Note for OPH Screen Microphysics}.
\href{https://github.com/FloatingPragma/observer-patch-holography/releases/download/\OPHPaperReleaseID/screen_microphysics_digital_calibration_note.pdf}{release PDF}.

\bibitem{oph-dark}
B.\ M\"uller,
\emph{Observer-Patch Holography and the Dark Matter Phenomenon}.
\href{https://github.com/FloatingPragma/observer-patch-holography/releases/download/\OPHPaperReleaseID/oph_dark_matter_paper.pdf}{release PDF}.

\bibitem{chinu-draft}
A.\ Osika and B.\ M\"uller,
\emph{\(\chinu\): Probing the OPH Dark Sector via Coherent-Matter Susceptibility},
draft, May 20, 2026.

\end{thebibliography}

\end{document}
